\section{Resultados e discussão}\label{sec:results}

A execução sintética considerou quatro verificações automáticas. Todas foram marcadas como aprovadas para demonstrar a forma de registro de resultados; esses valores não correspondem a um experimento científico externo e não devem ser reutilizados como evidência em outro trabalho.

\begin{table}[htbp]
\centering
\begin{tallabnttblr}
[
  caption={Resultados sintéticos da validação},
  label={tab:results},
  remark{Fonte}={Elaboração própria.},
  remark{Nota}={Dados exclusivamente demonstrativos.},
]
{
  colspec={X[l] X[c]},
}
\toprule
Verificação & Resultado \\
\midrule
Formato A4 & PASS \\
Fontes incorporadas & PASS \\
PDF/A & PASS \\
Marcadores estruturais & PASS \\
\bottomrule
\end{tallabnttblr}
\end{table}

Com $p=4$ e $n=4$, a Equação~\ref{eq:pass-rate} produz $R=100\%$. Esse valor deve ser interpretado apenas dentro do cenário construído. Um fluxo de validação real pode incluir dezenas de verificações, resultados de revisão e condições específicas do perfil documental.

\subsection{Interpretação}

Os resultados ilustram duas vantagens da automação. A primeira é a possibilidade de detectar problemas repetitivos antes da revisão final. A segunda é registrar o mesmo procedimento em execuções futuras, reduzindo a dependência de conferências manuais para propriedades objetivas.

Por outro lado, um conjunto de testes aprovado não prova que o texto está academicamente correto. Um PDF pode apresentar dimensões, fontes e referências tecnicamente válidas e ainda conter problemas de argumentação, metodologia ou interpretação. Assim, a automação deve complementar, e não substituir, a revisão do autor, do orientador e das instâncias institucionais aplicáveis.

\subsection{Uso do modelo}

No contexto deste TCC tutorial, o principal resultado é a própria organização do projeto: configuração concentrada no arquivo principal, capítulos separados, bibliografia estruturada e objetos acadêmicos inseridos onde são efetivamente discutidos. Essa organização permite que o estudante substitua o conteúdo demonstrativo pelo seu trabalho sem precisar conservar páginas de documentação dentro do PDF final.
